% Irodov, section 2.4 problems 122--129, printed pp.76--77,
% PDF pp.85--86. Fig 2.3 original inspected for #127.
\thermoentropyproblem{122}{76}
理想气体经历由\emph{一条等温线、
一条多方线及一条绝热线}组成的可逆热机循环，
等温阶段位于循环的最高温度 $T_h$。
求热效率，并以循环最高温度与最低温度之比
$\tau=T_h/T_c>1$ 表示。
\exsolution
设多方指数为 $n$，与绝热指数 $\gamma$ 不同，
并取其支路在冷却时热容量为正，
使循环沿正功方向运行。
沿两条非等温支路，
可逆绝热满足
$TV^{\gamma-1}=\text{常数}$，
多方满足
$TV^{n-1}=\text{常数}$。
设高温等温膨胀的体积倍率为 $r$。
由两个连接段分别跨越相同的温比 $\tau$，
\[
\ln r=
\ln\tau\left(\frac1{\gamma-1}-\frac1{n-1}\right).
\]
在高温等温支路吸热
\[
Q_{\rm in}=\nu RT_h\ln r
=\nu RT_h\ln\tau
\left(\frac1{\gamma-1}-\frac1{n-1}\right).
\]
多方冷却支路的摩尔热容为
\[
C_n=C_V+\frac R{1-n}
=R\left(\frac1{\gamma-1}-\frac1{n-1}\right).
\]
它的排热量为
$Q_{\rm out}=\nu C_n(T_h-T_c)$，
绝热支路无热交换。
因此
\[
\frac{Q_{\rm out}}{Q_{\rm in}}
=\frac{T_h-T_c}{T_h\ln\tau}
=\frac{\tau-1}{\tau\ln\tau},
\]
\[
\boxed{\eta_{\rm cycle}
=1-\frac{\tau-1}{\tau\ln\tau}}.
\]
\excheck 本版第122题题面未给出温度倍率的数值，
所以没有可计算的数值效率。
这里显式定义 $\tau=T_h/T_c$，
并假设热源过程方向与可逆热机一致；
与第121题的函数形式相同，
不是给题目虚构了一个数值倍率。

\thermoentropyproblem{123}{77}
一理想气体的顺时针正功循环
由一条可逆绝热、一条等压和一条等容线组成。
气体绝热指数为 $\gamma$。
如果绝热阶段体积
(a) 增大 $n$ 倍；
(b) 缩小 $n$ 倍，
求对应的热效率。
\exsolution
(a) 令绝热膨胀初态为
$A:(p_1,V_1,T_1)$，
末态
$B:(p_1n^{-\gamma},nV_1,T_1n^{1-\gamma})$。
然后在低压下等压压缩至
$C:(p_1n^{-\gamma},V_1,T_1n^{-\gamma})$，
最后等容加热回 $A$。
吸热发生于 $C\to A$，
排热发生于 $B\to C$：
\[
Q_{\rm in}=\nu C_VT_1(1-n^{-\gamma}),
\quad
Q_{\rm out}=\nu C_pT_1n^{-\gamma}(n-1).
\]
所以
\[
\boxed{\eta_a=
1-\frac{\gamma(n-1)}{n^\gamma-1}}.
\]
(b) 令绝热压缩初态为
$A:(p_1,nV_1,T_1)$，
压缩至
$B:(p_1n^\gamma,V_1,T_1n^{\gamma-1})$。
然后在高压下等压膨胀至
$C:(p_1n^\gamma,nV_1,T_1n^\gamma)$，
最后等容冷却回 $A$。
吸热在 $B\to C$，
排热在 $C\to A$：
\[
Q_{\rm in}=\nu C_pT_1n^{\gamma-1}(n-1),
\quad Q_{\rm out}=\nu C_VT_1(n^\gamma-1).
\]
故
\[
\boxed{\eta_b=
1-\frac{n^\gamma-1}
{\gamma n^{\gamma-1}(n-1)}}.
\]
\excheck 必须明确绝热阶段是膨胀还是压缩，
否则会混淆两条正功循环的吸放热支路。

\thermoentropyproblem{124}{77}
由等温、等压及等容三段组成的理想气体
顺时针热机循环，绝热指数为 $\gamma$。
若等温支路上的体积
(a) 增大 $n$ 倍；(b) 缩小 $n$ 倍，
求循环效率。
\exsolution
(a) 从状态
$A:(p_1,V_1,T_1)$ 等温膨胀到
$B:(p_1/n,nV_1,T_1)$。
然后在低压 $p_1/n$ 下等压压缩至
$C:(p_1/n,V_1,T_1/n)$，
最后等容加热回 $A$。
吸热来自高温等温与末段等容加热：
\[
Q_{\rm in}=\nu RT_1\ln n+
\nu C_VT_1(1-1/n),
\]
排热来自等压压缩：
$Q_{\rm out}=\nu C_pT_1(1-1/n)$。
计算效率：
\[
\boxed{\eta_a=
1-\frac{\gamma(n-1)}
{n-1+(\gamma-1)n\ln n}}.
\]
(b) 从
$A:(p_1/n,nV_1,T_1)$
等温压缩到
$B:(p_1,V_1,T_1)$。
然后在高压 $p_1$ 下等压膨胀至
$C:(p_1,nV_1,nT_1)$，
最后等容冷却回 $A$。
吸热仅在等压膨胀：
$Q_{\rm in}=\nu C_pT_1(n-1)$；
排热为等温压缩的
$\nu RT_1\ln n$ 和等容冷却的
$\nu C_VT_1(n-1)$。
所以
\[
\boxed{\eta_b=
1-\frac{n-1+(\gamma-1)\ln n}
{\gamma(n-1)}}.
\]
\excheck (a) 的热输入有\emph{两条支路}，
而 (b) 的热输出有两条支路，
应逐项统计而非只套普通卡诺效率。

\thermoentropyproblem{125}{77}
由两条等容线和两条等温线构成理想气体
正功循环，工质绝热指数为 $\gamma$。
整个循环最大体积与最小体积之比为
$\upsilon>1$，最高温与最低温之比为
$\tau>1$。求\emph{未配置回热器}的循环效率。
\exsolution
取高温等温膨胀段温度 $T_h$，
低温等温压缩段温度 $T_c=T_h/\tau$，
体积倍率
$\upsilon=V_{\max}/V_{\min}$。
循环净功等于两个等温功之差：
\[
W=\nu R(T_h-T_c)\ln\upsilon.
\]
吸热不仅来自高温等温膨胀，
还来自在最小体积处的等容加热：
\[
Q_{\rm in}=\nu RT_h\ln\upsilon+
\nu C_V(T_h-T_c).
\]
所以
\[
\eta=
\frac{(T_h-T_c)\ln\upsilon}
{T_h\ln\upsilon+(C_V/R)(T_h-T_c)}.
\]
用 $C_V/R=1/(\gamma-1)$ 及
$T_h/T_c=\tau$：
\[
\boxed{\eta=
\frac{(\tau-1)\ln\upsilon}
{\tau\ln\upsilon+(\tau-1)/(\gamma-1)}}.
\]
\excheck 若理想回热器完全回收等容冷却的热，
外部吸热计算不同；原书本题未设回热器，
所以必须保留等容加热项。

\thermoentropyproblem{126}{77}
理想气体经由两条等压线和两条等温线组成
顺时针热机循环，绝热指数为 $\gamma$。
循环最高和最低压强之比为 $n>1$，
最高与最低温之比为 $\tau>1$。
求\emph{无回热}循环的效率。
\exsolution
高温等温膨胀的压强从高到低，
低温等温压缩的压强从低到高，
等温压强比均为 $n$，
故循环净功
\[
W=\nu R(T_h-T_c)\ln n.
\]
外界向气体供热在
高温等温膨胀和高压等压加热两段：
\[
Q_{\rm in}=\nu RT_h\ln n+
\nu C_p(T_h-T_c).
\]
于是
\[
\eta=\frac{(T_h-T_c)\ln n}
{T_h\ln n+(C_p/R)(T_h-T_c)}.
\]
使用
$C_p/R=\gamma/(\gamma-1)$，
\[
\boxed{\eta=
\frac{(\tau-1)\ln n}
{\tau\ln n+\gamma(\tau-1)/(\gamma-1)}}.
\]
\excheck 热量计入外界供热，
不假设存在理想回热器；
若忽略高压等压加热，效率将被高估。

\thermoentropyproblem{127}{77}
原书图2.3的顺时针 $p$–$V$ 三角形循环中，
一边为过原点的直线 $p=\alpha V$，
其余两边分别为等容线与等压线。
若循环最高温和最低温之比为
$\tau>1$，工质的绝热指数为 $\gamma$，
求效率。
\exsolution
取最低温点
$A=(V_1,p_1)$，最高温点
$B=(V_2,p_2)$，
二者在直线 $p=\alpha V$ 上，
因此若 $\lambda=V_2/V_1>1$，
则 $p_2/p_1=\lambda$，
$T_B/T_A=\lambda^2$，
故 $\lambda=\sqrt\tau$。
按原图方向，
$A\to B$ 沿 $p\propto V$ 吸热，
$B\to C$ 在 $V_2$ 等容冷却，
$C\to A$ 在 $p_1$ 等压压缩。
沿 $p\propto V$ 的摩尔热容为
$C_{AB}=R(\gamma+1)/[2(\gamma-1)]$，
于是
\[
Q_{\rm in}
=\nu\frac{R(\gamma+1)}{2(\gamma-1)}
(T_B-T_A)
=\frac{\gamma+1}{2(\gamma-1)}
p_1V_1(\lambda^2-1).
\]
闭合循环净功为三角形面积：
\[
W=\frac12(p_2-p_1)(V_2-V_1)
=\frac12p_1V_1(\lambda-1)^2.
\]
相除：
\[
\boxed{\eta=
\frac{\gamma-1}{\gamma+1}
\frac{\sqrt\tau-1}{\sqrt\tau+1}}.
\]
\excheck 原图箭头是\emph{顺时针}；
若颠倒运动方向则为需输入外功的逆循环，
不能把负功解释成热机效率。

\thermoentropyproblem{128}{77}
利用 Clausius 不等式证明：
在具有相同最高温度 $T_{\max}$ 和
最低温度 $T_{\min}$ 的热机循环中，
没有任何循环能比在这两个温度间工作的
可逆卡诺循环效率更高。
\exsolution
令某循环吸热段从环境吸收总热量 $Q_{\rm in}>0$，
放热段向环境放出热量模 $Q_{\rm out}>0$。
Clausius 不等式
\[
\oint\frac{\delta Q}{T}\le0
\]
可拆为吸热与放热两部分。
因为所有吸热发生的温度
$T\le T_{\max}$，
所以
$\int_{\rm in}\delta Q/T\ge Q_{\rm in}/T_{\max}$。
所有放热发生的温度
$T\ge T_{\min}$，
所以
$\int_{\rm out}\delta Q_{\rm out}/T
\le Q_{\rm out}/T_{\min}$。
因此
\[
0\ge
\oint\frac{\delta Q}{T}
=\int_{\rm in}\frac{\delta Q}{T}
-\int_{\rm out}\frac{\delta Q_{\rm out}}T
\ge\frac{Q_{\rm in}}{T_{\max}}
-\frac{Q_{\rm out}}{T_{\min}}.
\]
于是
$Q_{\rm out}/Q_{\rm in}\ge T_{\min}/T_{\max}$，
所以
\[
\boxed{\eta=\frac{W}{Q_{\rm in}}
=1-\frac{Q_{\rm out}}{Q_{\rm in}}
\le1-\frac{T_{\min}}{T_{\max}}
=\eta_{\rm Carnot}}.
\]
\excheck 等号要求理想可逆且吸热全在
$T_{\max}$、放热全在 $T_{\min}$
进行；一般不可逆循环严格低于此上界。

\thermoentropyproblem{129}{77}
利用卡诺定理证明：对由温度 $T$ 和体积 $V$
决定状态的均匀简单可压缩物质，
其内能 $U(T,V)$ 与状态方程 $p(T,V)$ 满足
\[
\left(\frac{\partial U}{\partial V}\right)_T
=T\left(\frac{\partial p}{\partial T}\right)_V-p.
\]
提示：考虑 $p$–$V$ 图上的无穷小卡诺循环。
\exsolution
在温度 $T$ 和 $T+\mathrm dT$ 之间
做无穷小可逆卡诺循环，
高温等温阶段的体积变化为 $\mathrm dV$。
卡诺定理给出效率的一阶式
\[
\frac{\delta W}{\delta Q_h}=\frac{\mathrm dT}{T}.
\]
在无穷小循环极限，
等温线之间的压差为
$\mathrm dp=(\partial p/\partial T)_V\,\mathrm dT$，
故循环净功的一阶项为
\[
\delta W=\left(\frac{\partial p}{\partial T}\right)_V
\mathrm dT\,\mathrm dV.
\]
高温等温支路吸热由第一定律
\[
\delta Q_h=\mathrm dU+p\,\mathrm dV
=\left[\left(\frac{\partial U}{\partial V}\right)_T+p\right]
\mathrm dV.
\]
代入卡诺效率并约去微元：
\[
\frac{(\partial p/\partial T)_V\,\mathrm dT}
{(\partial U/\partial V)_T+p}
=\frac{\mathrm dT}{T}.
\]
因此
\[
\boxed{\left(\frac{\partial U}{\partial V}\right)_T
=T\left(\frac{\partial p}{\partial T}\right)_V-p}.
\]
\excheck 一阶近似忽略的项是比
$\mathrm dT\,\mathrm dV$ 更高阶的微小面积贡献；
所得结果也可由 Helmholtz 自由能的
Maxwell 关系独立复核。
